\documentclass[10pt,a4paper]{amsart}
\usepackage{amsmath,amssymb,amsthm,graphicx}
% Retained minimal preamble for the Stacks Project chapter coherent.tex (Stacks chapter, pre-delegation retained fragment for stacks-coherent).
% Verbatim excerpt of the chapter preamble (sources/stacks-spaces-descent/preamble.tex): the reference/slogan/history
% environment block (native lines 7-12), Xy-pic setup (native lines 14-19) and the native theorem declarations with
% the mathematical macros (native lines 154-217). Copied mechanically; nothing added, removed or redefined. The
% wrapper class and the amsmath/amssymb/amsthm/graphicx base packages are supplied by the assembler, so the
% class/xr-hyper/hyperref/externaldocument part of the native preamble is omitted.
% For dealing with references we use the comment environment
\usepackage{verbatim}
\newenvironment{reference}{\comment}{\endcomment}
%\newenvironment{reference}{}{}
\newenvironment{slogan}{\comment}{\endcomment}
\newenvironment{history}{\comment}{\endcomment}
% For commutative diagrams we use Xy-pic
\usepackage[all]{xy}

% We use 2cell for 2-commutative diagrams.
\xyoption{2cell}
\UseAllTwocells
% Theorem environments.
%
\theoremstyle{plain}
\newtheorem{theorem}[subsection]{Theorem}
\newtheorem{proposition}[subsection]{Proposition}
\newtheorem{lemma}[subsection]{Lemma}

\theoremstyle{definition}
\newtheorem{definition}[subsection]{Definition}
\newtheorem{example}[subsection]{Example}
\newtheorem{exercise}[subsection]{Exercise}
\newtheorem{situation}[subsection]{Situation}

\theoremstyle{remark}
\newtheorem{remark}[subsection]{Remark}
\newtheorem{remarks}[subsection]{Remarks}

\numberwithin{equation}{subsection}

% Macros
%
\def\lim{\mathop{\mathrm{lim}}\nolimits}
\def\colim{\mathop{\mathrm{colim}}\nolimits}
\def\Spec{\mathop{\mathrm{Spec}}}
\def\Hom{\mathop{\mathrm{Hom}}\nolimits}
\def\Ext{\mathop{\mathrm{Ext}}\nolimits}
\def\SheafHom{\mathop{\mathcal{H}\!\mathit{om}}\nolimits}
\def\SheafExt{\mathop{\mathcal{E}\!\mathit{xt}}\nolimits}
\def\Sch{\mathit{Sch}}
\def\Mor{\mathop{\mathrm{Mor}}\nolimits}
\def\Ob{\mathop{\mathrm{Ob}}\nolimits}
\def\Sh{\mathop{\mathit{Sh}}\nolimits}
\def\NL{\mathop{N\!L}\nolimits}
\def\CH{\mathop{\mathrm{CH}}\nolimits}
\def\proetale{{pro\text{-}\acute{e}tale}}
\def\etale{{\acute{e}tale}}
\def\QCoh{\mathit{QCoh}}
\def\Ker{\mathop{\mathrm{Ker}}}
\def\Im{\mathop{\mathrm{Im}}}
\def\Coker{\mathop{\mathrm{Coker}}}
\def\Coim{\mathop{\mathrm{Coim}}}

% Boxtimes
%
\DeclareMathSymbol{\boxtimes}{\mathbin}{AMSa}{"02}

%
% Macros for moduli stacks/spaces
%
\def\QCohstack{\mathcal{QC}\!\mathit{oh}}
\def\Cohstack{\mathcal{C}\!\mathit{oh}}
\def\Spacesstack{\mathcal{S}\!\mathit{paces}}
\def\Quotfunctor{\mathrm{Quot}}
\def\Hilbfunctor{\mathrm{Hilb}}
\def\Curvesstack{\mathcal{C}\!\mathit{urves}}
\def\Polarizedstack{\mathcal{P}\!\mathit{olarized}}
\def\Complexesstack{\mathcal{C}\!\mathit{omplexes}}
% \Pic is the operator that assigns to X its picard group, usage \Pic(X)
% \Picardstack_{X/B} denotes the Picard stack of X over B
% \Picardfunctor_{X/B} denotes the Picard functor of X over B
\def\Pic{\mathop{\mathrm{Pic}}\nolimits}
\def\Picardstack{\mathcal{P}\!\mathit{ic}}
\def\Picardfunctor{\mathrm{Pic}}
\def\Deformationcategory{\mathcal{D}\!\mathit{ef}}
\usepackage{cleveref}
\crefname{originallemma}{original lemma}{original lemmas}
\Crefname{originallemma}{Original lemma}{Original lemmas}
\makeatletter
\newlabel{cohomology-lemma-projection-formula}{{01E8}{source}{}{sourcepointer-1}{}}
\newlabel{cohomology-lemma-projection-formula@cref}{{[originallemma][0][]01E8}{[1][0][]source}{}{}{}}
\makeatother
\crefname{originallemma}{original lemma}{original lemmas}
\Crefname{originallemma}{Original lemma}{Original lemmas}
\makeatletter
\newlabel{derived-lemma-two-ss-complex-functor}{{015J}{source}{}{sourcepointer-2}{}}
\newlabel{derived-lemma-two-ss-complex-functor@cref}{{[originallemma][0][]015J}{[1][0][]source}{}{}{}}
\makeatother
\crefname{originallemma}{original lemma}{original lemmas}
\Crefname{originallemma}{Original lemma}{Original lemmas}
\makeatletter
\newlabel{lemma-cohomology-projective-space-over-base}{{01XW}{source}{}{sourcepointer-3}{}}
\newlabel{lemma-cohomology-projective-space-over-base@cref}{{[originallemma][0][]01XW}{[1][0][]source}{}{}{}}
\makeatother
\crefname{originallemma}{original lemma}{original lemmas}
\Crefname{originallemma}{Original lemma}{Original lemmas}
\makeatletter
\newlabel{lemma-cohomology-projective-space-over-ring}{{01XT}{source}{}{sourcepointer-4}{}}
\newlabel{lemma-cohomology-projective-space-over-ring@cref}{{[originallemma][0][]01XT}{[1][0][]source}{}{}{}}
\makeatother
\crefname{originallemma}{original lemma}{original lemmas}
\Crefname{originallemma}{Original lemma}{Original lemmas}
\makeatletter
\newlabel{lemma-identify-functorially}{{01XV}{source}{}{sourcepointer-5}{}}
\newlabel{lemma-identify-functorially@cref}{{[originallemma][0][]01XV}{[1][0][]source}{}{}{}}
\makeatother
\crefname{originallemma}{original lemma}{original lemmas}
\Crefname{originallemma}{Original lemma}{Original lemmas}
\makeatletter
\newlabel{more-algebra-lemma-homotopy-koszul-abstract}{{0626}{source}{}{sourcepointer-6}{}}
\newlabel{more-algebra-lemma-homotopy-koszul-abstract@cref}{{[originallemma][0][]0626}{[1][0][]source}{}{}{}}
\makeatother
\title{Cohomology of a projective bundle: for S a scheme, E a finite locally free O\_S-module of constant rank n + 1 and pi : P(E) -> S, the higher direct images R\textasciicircum{}q pi\_* O\_\{P(E)\}(d) equal Sym\textasciicircum{}d(E) for q = 0, 0 for q not in \{0, n\}, and SheafHom(Sym\textasciicircum{}\{-n-1-d\}(E) tensor wedge\textasciicircum{}\{n+1\} E, O\_S) for q = n; the identifications are compatible with base change and isomorphisms of locally free sheaves.}
\author{The Stacks Project Authors}
\date{stacks-project:coherent.tex:local-snapshot-20261012; GNU FDL 1.2 or later}
\begin{document}
\maketitle
\noindent\textit{Editorial note.} Statement span is lines 2271-2298 of sources/stacks-coherent/pinned-original-source.tex; proof span is lines 2300-2412, the complete author proof (111 body lines).\par\medskip
\noindent\textit{Editorial note.} Navigation stubs: 6 cited-but-not-included labels are supplied as non-load-bearing pointers with their pinned tags where the snapshot tag index has them.\par\medskip
\noindent\textit{Editorial note.} extra\_packages is empty on purpose: the retained preamble already loads xy with options and a second declaration clashes in the assembler wrapper.\par\medskip
\noindent\textit{Editorial note.} Text is retained verbatim, including the author's own typos and placeholders.\par\medskip
\noindent\textit{Editorial note.} Pool stacks-coherent (118 LIVE-clean, CONTROL.md section 2 ruling 14:3x). Top-20 triage: 9 candidates matched public corpus propositions (PUBLISHED, rejected with trace); 11 accepted here; next 9 pool ranks backfill.\par\medskip
\noindent\textit{Editorial note.} Span computed from the \textbackslash{}label marker through the enclosing \textbackslash{}end\{lemma\}; precard\_gate combined span RESULT=OK (probe L9999).\par\medskip
\noindent\textit{Allowed standard prerequisites (editorial).}\begin{itemize}
\item Explicit prerequisite boundary. See why-field; every pointer named here occurs verbatim in the retained statement/proof text.
\end{itemize}\medskip
\ifdefined\hypertarget\hypertarget{sourcepointer-1}{}\fi
\noindent\textit{Source reference.} Original lemma 01E8: sources/stacks-coherent/pinned-original-source.tex, source lines 2412--2412. Cited by the retained statement or proof but not part of the collected spans; supplied only as a navigation pointer. The argument does not depend on its content.\par\medskip
\ifdefined\hypertarget\hypertarget{sourcepointer-2}{}\fi
\noindent\textit{Source reference.} Original lemma 015J: sources/stacks-coherent/pinned-original-source.tex, source lines 2412--2412. Cited by the retained statement or proof but not part of the collected spans; supplied only as a navigation pointer. The argument does not depend on its content.\par\medskip
\ifdefined\hypertarget\hypertarget{sourcepointer-3}{}\fi
\noindent\textit{Source reference.} Original lemma 01XW: sources/stacks-coherent/pinned-original-source.tex, source lines 2222--2239. Cited by the retained statement or proof but not part of the collected spans; supplied only as a navigation pointer. The argument does not depend on its content.\par\medskip
\ifdefined\hypertarget\hypertarget{sourcepointer-4}{}\fi
\noindent\textit{Source reference.} Original lemma 01XT: sources/stacks-coherent/pinned-original-source.tex, source lines 1820--1823. Cited by the retained statement or proof but not part of the collected spans; supplied only as a navigation pointer. The argument does not depend on its content.\par\medskip
\ifdefined\hypertarget\hypertarget{sourcepointer-5}{}\fi
\noindent\textit{Source reference.} Original lemma 01XV: sources/stacks-coherent/pinned-original-source.tex, source lines 2102--2121. Cited by the retained statement or proof but not part of the collected spans; supplied only as a navigation pointer. The argument does not depend on its content.\par\medskip
\ifdefined\hypertarget\hypertarget{sourcepointer-6}{}\fi
\noindent\textit{Source reference.} Original lemma 0626: sources/stacks-coherent/pinned-original-source.tex, source lines 2412--2412. Cited by the retained statement or proof but not part of the collected spans; supplied only as a navigation pointer. The argument does not depend on its content.\par\medskip
\begin{lemma}
\label{lemma-cohomology-projective-bundle}
Let $S$ be a scheme. Let $n \geq 1$.
Let $\mathcal{E}$ be a finite locally
free $\mathcal{O}_S$-module of constant rank $n + 1$.
Consider the structure morphism
$$
\pi : \mathbf{P}(\mathcal{E}) \longrightarrow S.
$$
We have
$$
R^q\pi_*(\mathcal{O}_{\mathbf{P}(\mathcal{E})}(d)) =
\left\{
\begin{matrix}
\text{Sym}^d(\mathcal{E}) & \text{if} & q = 0 \\
0 & \text{if} & q \not = 0, n \\
\SheafHom_{\mathcal{O}_S}(
\text{Sym}^{- n - 1 - d}(\mathcal{E})
\otimes_{\mathcal{O}_S}
\wedge^{n + 1}\mathcal{E},
\mathcal{O}_S)
& \text{if} & q = n
\end{matrix}
\right.
$$
These identifications are compatible with base change and
isomorphism between locally free sheaves.
\end{lemma}

\begin{proof}
Consider the canonical map
$$
\pi^*\mathcal{E} \longrightarrow \mathcal{O}_{\mathbf{P}(\mathcal{E})}(1)
$$
and twist down by $1$ to get
$$
\pi^*(\mathcal{E})(-1) \longrightarrow \mathcal{O}_{\mathbf{P}(\mathcal{E})}
$$
This is a surjective map from a locally free rank $n + 1$ sheaf onto
the structure sheaf. Hence the corresponding Koszul complex is
exact (More on Algebra, Lemma
\ref{more-algebra-lemma-homotopy-koszul-abstract}).
In other words there is an exact complex
$$
0 \to
\pi^*(\wedge^{n + 1}\mathcal{E})(-n - 1) \to
\ldots \to
\pi^*(\wedge^i\mathcal{E})(-i) \to
\ldots \to
\pi^*\mathcal{E}(-1) \to
\mathcal{O}_{\mathbf{P}(\mathcal{E})} \to 0
$$
We will think of the term $\pi^*(\wedge^i\mathcal{E})(-i)$ as being
in degree $-i$.
We are going to compute the higher direct images
of this acyclic complex using the first spectral sequence of
Derived Categories, Lemma \ref{derived-lemma-two-ss-complex-functor}.
Namely, we see that there is a spectral sequence with terms
$$
E_1^{p, q} = R^q\pi_*\left(\pi^*(\wedge^{-p}\mathcal{E})(p)\right)
$$
converging to zero! By the projection formula
(Cohomology, Lemma \ref{cohomology-lemma-projection-formula})
we have
$$
E_1^{p, q} = \wedge^{-p} \mathcal{E} \otimes_{\mathcal{O}_S}
R^q\pi_*\left(\mathcal{O}_{\mathbf{P}(\mathcal{E})}(p)\right).
$$
Note that locally on $S$ the sheaf $\mathcal{E}$ is trivial,
i.e., isomorphic to $\mathcal{O}_S^{\oplus n + 1}$, hence locally on
$S$ the morphism $\mathbf{P}(\mathcal{E}) \to S$ can be identified
with $\mathbf{P}^n_S \to S$. Hence
locally on $S$ we can use the result of Lemmas
\ref{lemma-cohomology-projective-space-over-ring},
\ref{lemma-identify-functorially}, or
\ref{lemma-cohomology-projective-space-over-base}.
It follows that $E_1^{p, q} = 0$ unless $(p, q)$ is $(0, 0)$
or $(-n - 1, n)$. The nonzero terms are
\begin{align*}
E_1^{0, 0} & = \pi_*\mathcal{O}_{\mathbf{P}(\mathcal{E})} = \mathcal{O}_S \\
E_1^{-n - 1, n} & =
R^n\pi_*\left(\pi^*(\wedge^{n + 1}\mathcal{E})(-n - 1)\right) =
\wedge^{n + 1}\mathcal{E} \otimes_{\mathcal{O}_S}
R^n\pi_*\left(\mathcal{O}_{\mathbf{P}(\mathcal{E})}(-n - 1)\right)
\end{align*}
Hence there can only be one nonzero
differential in the spectral sequence namely the map
$d_{n + 1}^{-n - 1, n} : E_{n + 1}^{-n - 1, n} \to E_{n + 1}^{0, 0}$
which has to be an isomorphism (because the spectral sequence converges
to the $0$ sheaf). Thus $E_1^{p, q} = E_{n + 1}^{p, q}$ and
we obtain a canonical isomorphism
$$
\wedge^{n + 1}\mathcal{E} \otimes_{\mathcal{O}_S}
R^n\pi_*\left(\mathcal{O}_{\mathbf{P}(\mathcal{E})}(-n - 1)\right) =
R^n\pi_*\left(\pi^*(\wedge^{n + 1}\mathcal{E})(-n - 1)\right)
\xrightarrow{d_{n + 1}^{-n - 1, n}}
\mathcal{O}_S
$$
Since $\wedge^{n + 1}\mathcal{E}$ is an invertible
sheaf, this implies that
$R^n\pi_*\mathcal{O}_{\mathbf{P}(\mathcal{E})}(-n - 1)$ is invertible
as well and canonically isomorphic to the inverse of
$\wedge^{n + 1}\mathcal{E}$. In other words we have proved the case
$d = - n - 1$ of the lemma.

\medskip\noindent
Working locally on $S$ we see immediately from the computation of
cohomology in Lemmas \ref{lemma-cohomology-projective-space-over-ring},
\ref{lemma-identify-functorially}, or
\ref{lemma-cohomology-projective-space-over-base} the statements on
vanishing of the lemma. Moreover the result on $R^0\pi_*$ is clear
as well, since there are canonical maps
$\text{Sym}^d(\mathcal{E}) \to \pi_* \mathcal{O}_{\mathbf{P}(\mathcal{E})}(d)$
for all $d$. It remains to show that the description of
$R^n\pi_*\mathcal{O}_{\mathbf{P}(\mathcal{E})}(d)$ is correct
for $d < -n - 1$. In order to do this we consider the map
$$
\pi^*(\text{Sym}^{-d - n - 1}(\mathcal{E}))
\otimes_{\mathcal{O}_{\mathbf{P}(\mathcal{E})}}
\mathcal{O}_{\mathbf{P}(\mathcal{E})}(d)
\longrightarrow
\mathcal{O}_{\mathbf{P}(\mathcal{E})}(-n - 1)
$$
Applying $R^n\pi_*$ and the projection formula (see above) we get a map
$$
\text{Sym}^{-d - n - 1}(\mathcal{E})
\otimes_{\mathcal{O}_S}
R^n\pi_*(\mathcal{O}_{\mathbf{P}(\mathcal{E})}(d))
\longrightarrow
R^n\pi_*\mathcal{O}_{\mathbf{P}(\mathcal{E})}(-n - 1) =
(\wedge^{n + 1}\mathcal{E})^{\otimes -1}
$$
(the last equality we have shown above).
Again by the local calculations of Lemmas
\ref{lemma-cohomology-projective-space-over-ring},
\ref{lemma-identify-functorially}, or
\ref{lemma-cohomology-projective-space-over-base}
it follows that this map induces a perfect pairing between
$R^n\pi_*(\mathcal{O}_{\mathbf{P}(\mathcal{E})}(d))$ and
$\text{Sym}^{-d - n - 1}(\mathcal{E}) \otimes \wedge^{n + 1}(\mathcal{E})$
as desired.
\end{proof}

\end{document}
